problem:  
Let \( M^n \subset (S^{n+1},g) \) be a closed, two-sided minimal hypersurface, where \(g\) is a Riemannian metric \(C^2\)-close to the round metric and satisfies the sectional curvature pinching condition  
\[
(1-\varepsilon)\leq K_g \leq (1+\varepsilon)
\]
for some small \(\varepsilon>0\). Denote by \(\mathrm{Ind}_g(M)\) the Morse index of \(M\) with respect to the area functional under metric \(g\).  

For the round metric (\(\varepsilon=0\)), the index behavior of non-totally geodesic minimal hypersurfaces is partially understood (e.g. the Clifford hypersurface in \(S^{n+1}\) gives \(\mathrm{Ind}=n+3\)).  

**Research question:**  
Does there exist a universal constant \(c_n>0\) such that, for every sufficiently small \(\varepsilon>0\) and every non-totally-geodesic minimal hypersurface \(M^n \subset (S^{n+1},g)\) as above, one has  
\[
\mathrm{Ind}_g(M) \ge c_n\, (1 - C\varepsilon)
\]
for some constant \(C=C(n)\) independent of \(M\)?  
Equivalently, is the “minimal index-gap rigidity" of the round sphere stable under \(C^2\)-small metric perturbations with positive Ricci curvature?  

In other words, can the near-round geometry of \((S^{n+1},g)\) be detected from lower bounds on the Morse index spectrum of minimal hypersurfaces that are stable under such curvature pinching?

Why is it a "good" problem:  
This revised question imposes precise geometric control — \(C^2\)-small perturbations and sectional curvature pinching — making the relationship between ambient curvature and the stability operator of minimal hypersurfaces both well-defined and testable. The problem directly probes the stability of deep spectral–geometric phenomena: does the rigid spectral gap in the round case persist under small geometric deformations? It connects the theory of minimal hypersurfaces (index estimates, second variation geometry) with the study of metric perturbations and Ricci pinching in global differential geometry, thus bridging geometric analysis, perturbation theory, and spectral geometry.  
Its concise formulation — asking whether a uniform lower bound on index survives small curvature perturbations — hides substantial subtleties involving the analytic dependence of the Jacobi spectrum, curvature pinching, and rigidity versus flexibility of minimal hypersurfaces. This spectrum–rigidity linkage under perturbation is currently not understood, making the problem both sharply posed and potentially foundational for understanding stability phenomena in curved spaces.